\documentclass[11pt]{article}

\usepackage[T1]{fontenc}
\usepackage[utf8]{inputenc}
\usepackage{lmodern}
\usepackage{amsmath,amssymb,amsthm,mathtools}
\usepackage[a4paper,margin=28mm]{geometry}
\usepackage{microtype}
\usepackage{xurl}
\usepackage[hidelinks]{hyperref}
\usepackage{enumitem}
\usepackage{booktabs}
\usepackage{array}

\newtheorem{theorem}{Theorem}[section]
\newtheorem{lemma}[theorem]{Lemma}
\newtheorem{proposition}[theorem]{Proposition}
\newtheorem{corollary}[theorem]{Corollary}
\newtheorem{definition}[theorem]{Definition}
\newtheorem{remark}[theorem]{Remark}

\newcommand{\N}{\mathbb N}
\newcommand{\WIN}{\ensuremath{\mathrm{WIN}}}
\newcommand{\LOSS}{\ensuremath{\mathrm{LOSS}}}
\newcommand{\DRAW}{\ensuremath{\mathrm{DRAW}}}
\newcommand{\R}{\mathcal R}
\newcommand{\PhiM}{\Phi}
\newcommand{\odd}{\operatorname{odd}}
\newcommand{\vTwo}{v_2}
\newcommand{\source}{\operatorname{src}}

\title{The Two-Player $3n\mathbin{\pm}1$ Game:\\
A Binary Normal Form and a Well-Founded No-Draw Proof}
\author{Grisha Pochuev\\
  \small\href{mailto:n_854@mail.ru}{n\_854@mail.ru}\\
  \small\url{https://github.com/Grisha-Pochuev/3n-plus-minus-1-game}}
\date{Corrected manuscript, 10 August 2026}

\begin{document}
\maketitle

\begin{center}
\small\textbf{Audit status.} This manuscript presents the repaired human proof after
an external audit of an earlier version. The two invalid shortcuts identified by
that audit remain withdrawn. Independent external re-audit of the repaired bridge
is still pending, and no end-to-end Lean formalization is claimed.
\end{center}

\begin{abstract}
Consider the impartial two-player game on positive odd integers in which a move
from $n>1$ chooses $3n-1$ or $3n+1$ and removes all powers of two; reaching $1$
wins immediately. Arbitrary play need not terminate, since $5\to7\to5$ is a
legal cycle. The problem is whether optimal play nevertheless has finite
remoteness from every start. We pass to a conjugated binary game with an
expanding move $A$ and a strictly contracting move $B$, introduce constant-tail
coordinates, an embedded source map, and finite proof-height tokens for certified
winning and losing positions. The global argument is a well-founded marked
normalization: source decreases and certified token replacements form the outer
rank, while a typed factor/obligation system is well-founded inside each fixed
outer fibre. An external audit found that an earlier entry argument incorrectly
identified a smaller canonical coefficient with a smaller coefficient source and
also used a reverse-factor lemma outside its valid exponent range. We repair both
issues. The key final step is a one-shot entry component that permits a returned
inner source to be larger than the original source exactly once, during typed
normalizer initialization, while forbidding any repeated unranked reset. A direct
phase analysis of the remaining factor-free exponent-two base entry closes the
last attachment case. The resulting human proof rules out all draw positions.
Machine regressions and symbolic certificates are supporting checks with a
strictly narrower scope than the human theorem.
\end{abstract}

\tableofcontents

\section{The game and the theorem}

The current state is a positive odd integer $n$. For $n>1$ the player to move
chooses one of
\[
  3n-1,\qquad 3n+1,
\]
and divides the chosen value by $2$ until it is odd. If the resulting odd value
is $1$, the player who made the move wins. Players alternate with perfect
information.

For the player to move, a position is \WIN\ if some move leads to a \LOSS\
position, and is \LOSS\ if every move leads to a \WIN\ position. Because the
game graph is infinite, these inductively generated classes need not a priori
exhaust all positions. Every remaining position is called \DRAW. This is the
standard finite-remoteness interpretation of Conway's Beans-Don't-Talk game
\cite{Guy1986,Guy1996}; the modern prize formulation is given by Alth\"ofer
\cite{AlthoferPage,Althofer2024}.

\begin{theorem}[Main theorem]\label{thm:main}
Every positive odd starting position has finite remoteness. Equivalently, the
game has no \DRAW\ positions, and optimal play reaches $1$ after finitely many
moves.
\end{theorem}

The proof is not a termination proof for arbitrary play. For example,
\[
  5\longrightarrow7\longrightarrow5
\]
is a legal two-ply cycle under suitable choices.

\section{Binary conjugation}

Write an odd original state as $n=2m+1$ and put
\[
  F(m)=\left\lceil\frac{3m}{2}\right\rceil.
\]
For a nonnegative integer $x$, let $R(x)$ be the integer obtained by deleting
the maximal alternating suffix of the binary expansion of $x$. For example,
$R(1001_2)=10_2$ and $R(10101_2)=0$.

\begin{proposition}[First normal form]\label{prop:firstnormal}
For $m>0$, the two children in $m$-coordinates are exactly
\[
  F(m),\qquad F(R(m)).
\]
\end{proposition}

\begin{proof}
We have
\[
  3(2m+1)-1=2(3m+1),\qquad
  3(2m+1)+1=2(3m+2).
\]
Exactly one of $3m+1$ and $3m+2$ is odd. In the other expression, each
additional factor of two removed corresponds to crossing one more alternating
bit of the binary suffix of $m$. The division stops at the first repeated
adjacent bit, which is precisely deletion of the maximal alternating suffix.
The remaining odd value, returned to $m$-coordinates, is $F(R(m))$.
\end{proof}

Every child in Proposition~\ref{prop:firstnormal} lies in the image of $F$.
Representing $F(q)$ by $q$ gives the conjugated game
\[
  A(q)=F(q)=\left\lceil\frac{3q}{2}\right\rceil,
  \qquad B(q)=R(F(q)),
\]
with terminal state $q=0$.

\begin{proposition}[Strict contraction]\label{prop:Bcontracts}
For every $q>0$,
\[
  B(q)<q.
\]
\end{proposition}

\begin{proof}
Deleting a nonempty binary suffix gives
\[
 B(q)\le \left\lfloor\frac{F(q)}2\right\rfloor
 \le \left\lfloor\frac{3q+1}{4}\right\rfloor<q
\]
for $q\ge2$, and $B(1)=0$ directly.
\end{proof}

The difficulty is therefore concentrated in the expanding branch $A$ and in
the unbounded length of the alternating suffix deleted by $B$.

\section{Constant-tail coordinates and source coordinates}

For positive odd $a$, $r\ge1$, and $e\in\{0,1\}$ define
\[
  Q_r^e(a)=a2^r-e.
\]
This is the canonical form of a positive binary word with a maximal constant
suffix of $r$ copies of the bit $e$.

\begin{lemma}[Long-tail recurrence]\label{lem:longtail}
For every $r\ge3$,
\[
 \boxed{
 A(Q_r^e(a))=Q_{r-1}^e(3a),\qquad
 B(Q_r^e(a))=Q_{r-2}^e(3a).}
\]
Moreover the two boundary states share their contracting child:
\[
 \boxed{B(Q_1^e(a))=B(Q_2^e(a)).}
\]
\end{lemma}

\begin{proof}
For $r\ge3$, the expanding child still ends in at least two equal bits, so
its maximal alternating suffix consists only of the final bit. The boundary
identity follows by comparing the binary words of $3a-e$ and $6a-e$; the
additional appended bit extends the same alternating suffix and therefore
deletes to the same prefix.
\end{proof}

Define the embedded three-free coefficient
\[
  J(s)=2A(s)+1.
\]
Every positive odd coefficient has a unique factorization
\[
  a=3^kJ(s),\qquad k\ge0,
\]
which defines its coefficient source $s$.

\begin{lemma}[Embedded source transition]\label{lem:source-transition}
For $s>0$, the two signed boundary transitions from the coefficient $J(s)$
select exactly the two ordinary children $A(s)$ and $B(s)$. More precisely,
there is a phase
\[
 \alpha(s)=1-\left(\left\lfloor\frac s2\right\rfloor\bmod2\right)
\]
such that
\[
 \odd(3J(s)+1-2\alpha(s))=J(A(s))
\]
with $2$-adic valuation one, while the opposite phase selects $J(B(s))$ with
valuation at least two.
\end{lemma}

The point of this embedding is that multiplication of a constant-tail
coefficient by $3$ preserves its source, while a signed exponent-one boundary
reveals one ordinary $A/B$ move on the source.

\section{Finite outcome certificates and token rank}

Every certified \WIN\ or \LOSS\ position has a finite canonical proof height:
\[
 h(0)=0,
\]
\[
 h(x)=1+\min\{h(y):y\text{ is a \LOSS\ child of }x\}
 \quad(x\text{ \WIN}),
\]
and
\[
 h(x)=1+\max\{h(y):y\text{ is a child of }x\}
 \quad(x\text{ \LOSS}).
\]

A marked normalization configuration carries only finite states whose outcome
and height comparison have actually been proved. These states are called proof
tokens.

\begin{lemma}[Multiset token rank]\label{lem:multiset}
For a finite multiset $M$ of proof heights define
\[
 \PhiM(M)=\mathop{\#}_{h\in M}\omega^h,
\]
the natural ordinal sum of the monomials $\omega^h$. Replacing one token of
height $H$ by any finite multiset of certified descendants of heights strictly
below $H$ strictly decreases $\PhiM$.
\end{lemma}

\begin{proof}
In Cantor normal form the coefficient of $\omega^H$ decreases by one, while
all inserted terms have lower exponent. No collection of lower monomials can
compensate for the loss at exponent $H$.
\end{proof}

This rank permits an exceptional local diamond to split one token into several
certified descendants without losing well-foundedness.

\section{Universal factor scan}

Fix a three-free odd base $a=J(s)$ and a phase $e$. For $k\ge0$ put
\[
 P_k=Q_2^e(3^ka),\qquad R_k=Q_1^e(3^ka),
\]
\[
 C_k=B(P_k)=B(R_k),\qquad T_k=A(R_k).
\]
Then
\[
 \operatorname{moves}(P_k)=\{R_{k+1},C_k\},\qquad
 \operatorname{moves}(R_k)=\{T_k,C_k\}.
\]

\begin{lemma}[Two-level DRAW horizon]\label{lem:twolevel}
If the adjacent frame $\{P_k,R_k\}$ contains a \DRAW, then
\[
 \boxed{\{C_k,T_k,C_{k+1},T_{k+1}\}\text{ contains a \DRAW}.}
\]
\end{lemma}

\begin{proof}
The common child $C_k$ cannot be \LOSS, otherwise both frame members would be
\WIN. If $C_k$ is \DRAW\ we are done. If $C_k$ is \WIN\ and $R_k$ is \DRAW,
then $T_k$ is \DRAW. The only remaining case has $P_k$ \DRAW\ and $R_k$
non-draw; then $R_{k+1}$ is \DRAW, so one of its two children
$T_{k+1},C_{k+1}$ is \DRAW.
\end{proof}

Factor at either exposed level
\[
 3^{i+1}J(s)+1-2e=2^vJ(t),\qquad v\ge1.
\]
Then
\[
 T_i=Q_v^{1-e}(J(t)).
\]
If $v\ge2$, the raw sibling is exactly
\[
 \boxed{C_i=Q_{v-1}^{1-e}(J(t)).}
\]
If $v=1$ and the positive raw sibling is \DRAW, its coefficient source is
strictly smaller than $t$. Thus every factorful exponent-one \DRAW\ reaches a
factor-free lift, a factor-free adjacent frame, or a genuine source descent
within at most two consecutive factor levels.

A second reverse argument removes the unbounded factor exponent itself. If a
factor frame does not pull back to a factor-free \DRAW\ over the same source,
then its only residual orientation is a marked gate
\[
  P_k\text{ \LOSS},\qquad R_k\text{ \DRAW}.
\]
The exact predecessor $Q_2^e(3^{k-1}J(s))$ is then either \DRAW, in which case
factors of three are removed at exponent two, or \WIN, in which case its
other child is an actual \LOSS\ proof token. The nonexceptional side exposes a
strictly lower \WIN\ descendant; the exceptional side is an exact adjacent
frame over that \LOSS\ token. Hence no factor level creates an unmarked reset.

\section{The typed fixed-fibre normalizer}

The local source, factor, and proof-token lemmas form a finite collection of
routing macros. Their full case proofs are recorded in Sections 91--135 of the
proof supplement at the repository revision cited in
Section~\ref{sec:repro}. We use only the following assembled statement.

\begin{proposition}[Typed normalizer]\label{prop:typed-normalizer}
Suppose a marked entry has already been supplied with
\begin{enumerate}[label=(\roman*)]
\item a retained arithmetic source anchor,
\item a finite multiset of certified \WIN/\LOSS\ proof tokens with proved
      provenance, and
\item one of the typed obligation/factor control states of the supplement.
\end{enumerate}
Then the equal-retained-data routing relation is well-founded.
\end{proposition}

The crucial scope restriction is that temporary routing cursors are not rank
anchors. In particular a canonical continuation
\[
 O(x,e)\longrightarrow O(A(x),1-e)
\]
usually increases the displayed ordinary source $x$, but this $A(x)$ is only
